% ==================================================================
% Color Charge as Emergent Vertex Identity:
% The hDP Tetrahedral Skeleton and SU(3) as Bookkeeping in CPP
% Companion 15 to "Mechanistic Derivation of Relativistic Effects
% via Space Stress Vector (SSV) in the Dipole Sea" (Version 16)
% Version 1 — 19 March 2026
% ==================================================================

\documentclass[12pt]{article}

\usepackage{amsmath,amssymb,amsthm}
\usepackage{geometry}
\usepackage{hyperref}
\usepackage{booktabs}
\usepackage{enumitem}
\usepackage{parskip}
\usepackage{array}

\geometry{letterpaper, margin=1in}

\newtheorem{theorem}{Theorem}[section]
\newtheorem{proposition}[theorem]{Proposition}
\newtheorem{corollary}[theorem]{Corollary}
\newtheorem{remark}[theorem]{Remark}
\newtheorem{definition}[theorem]{Definition}

% ── CPP macros ────────────────────────────────────────────────────────────────
\newcommand{\lP}{l_P}
\newcommand{\mP}{m_P}
\newcommand{\PSR}{\mathrm{PSR}}
\newcommand{\SSV}{\mathrm{SSV}}
\newcommand{\LSP}{\mathrm{LSP}}
\newcommand{\eCP}{e\mathrm{CP}}
\newcommand{\qCP}{q\mathrm{CP}}
\newcommand{\eDP}{e\mathrm{DP}}
\newcommand{\qDP}{q\mathrm{DP}}
\newcommand{\hDP}{h\mathrm{DP}}
\newcommand{\ZDC}{\mathrm{ZDC}}
\newcommand{\ZBW}{\mathrm{ZBW}}
\newcommand{\as}{\alpha_s}
\newcommand{\Td}{T_d}
\newcommand{\Hf}{H_4}

\title{\textbf{Color Charge as Emergent Vertex Identity:\\
The hDP Tetrahedral Skeleton and SU(3)\\
as Bookkeeping in Conscious Point Physics}\\[6pt]
{\large Companion~15 to ``Mechanistic Derivation of Relativistic Effects\\
via Space Stress Vector (SSV) in the Dipole Sea'' (Version~16)}}

\author{Thomas Lee Abshier, ND\\
Hyperphysics Institute\\
\texttt{https://hyperphysics.com}\\
\texttt{drthomas007@protonmail.com}}

\date{19 March 2026 --- Version~1}

\begin{document}
\maketitle

\noindent\textbf{Keywords:} Conscious Point Physics, color charge, quark
confinement, hDP tetrahedron, SU(3), baryon structure, color neutrality,
vertex identity, emergent symmetry, DP Sea.

%======================================================================
\begin{abstract}
We derive the observed color charge properties of quarks from the
geometry of a spontaneously nucleated hybrid Dipole Particle ($\hDP$)
tetrahedral skeleton in the Dipole Sea, without postulating SU(3) as a
fundamental gauge symmetry.  The $\hDP$ tetrahedron --- a stable
alternating-charge skeleton of four hybrid DPs on 600-cell Grid Points
--- provides the bonding scaffold for baryons.  Three quarks occupy
three of its four vertices; the fourth vertex remains open, giving the
baryon its net charge.

The central theorem of this paper is: \emph{what QCD calls color charge
is vertex identity} --- a quark's ``color'' is simply which of the three
occupied tetrahedral vertices it is bonded to.  Color neutrality (why
hadrons are colorless) follows from complete three-vertex occupation.
The prohibition on free quarks and the restriction to specific hadron
configurations ($q\bar{q}$, $qqq$, $\bar{q}\bar{q}\bar{q}$) follow
from charge geometry.  SU(3) is identified as the permutation algebra of
3-of-4 vertex occupation --- a mathematical description of the
bookkeeping, not a fundamental physical input.

No new postulates beyond those of Version~16 are required.  The
correspondence between the 8 gluon modes and the 8 Gell-Mann generators
is noted and deferred to Companion~16.
\end{abstract}

%======================================================================
\section*{Plain Language Summary}
%======================================================================

In standard physics (QCD), quarks carry a property called ``color
charge'' --- red, green, or blue --- which determines how they bind
together.  The rules are: three quarks of different colors form a
proton or neutron; a quark and an antiquark of opposite colors form a
pion or other meson.  No free quarks are ever observed.  These rules
are encoded in the mathematics of SU(3), a gauge symmetry group.

CPP asks: what is the \emph{physical structure} underlying these
rules?  The answer is a tetrahedral skeleton.  The Dipole Sea
spontaneously nucleates stable tetrahedra built from hybrid Dipole
Particles ($\hDP$s) --- DPs that combine both electromagnetic and
strong-force character.  This skeleton has four vertices.  When three
quarks bond to three of those four vertices, we call their positions
``red,'' ``green,'' and ``blue.''  But these labels just mean
\emph{which vertex a quark is sitting on}.  Color charge is vertex
identity.

Color neutrality means all three quark-bonding vertices are occupied,
so no net strong charge is exposed.  A free quark would require a
chain extending to infinity (established in companion~14) --- which
is why quarks are never observed alone.  SU(3) is the mathematical
description of how many ways you can permute three objects on three
vertices --- useful bookkeeping, but not the underlying reality.

%======================================================================
\section{Introduction}
\label{sec:intro}
%======================================================================

Quantum Chromodynamics (QCD) describes the strong nuclear force through
SU(3) color gauge symmetry~\cite{gross1973,politzer1973}.  Quarks carry
one of three color charges (red, green, blue); gluons carry
color-anticolor combinations; hadrons are color-neutral bound states.
The mathematical structure is elegant and experimentally confirmed, but
SU(3) is a postulate: it is not derived from a more fundamental
physical substrate.

Companion~14~\cite{c14} derived quark confinement and the Cornell
potential from the qDP chaining mechanism of CPP, without invoking
SU(3).  The present companion addresses the complementary question:
\emph{what is the physical substrate of color charge itself?}

The answer is the $\hDP$ tetrahedral skeleton --- a spontaneously
nucleated geometric structure in the Dipole Sea that provides the
bonding scaffold for baryons.  Color charge, in this picture, is not
an intrinsic quantum number carried by a quark.  It is a \emph{vertex
label}: red, green, and blue are names for the three occupied bonding
sites on the tetrahedral skeleton.

This reframing has several consequences:
\begin{enumerate}[noitemsep]
\item Color neutrality is a geometric condition, not a gauge condition.
\item The prohibition on free quarks is the topological condition
  derived in companion~14, now given a structural complement.
\item SU(3) is identified as the permutation algebra of the geometry
  --- correct, useful, but not fundamental.
\item The 3-of-4 vertex occupation rule explains the net charge of
  baryons without a separate assumption.
\end{enumerate}

Section~\ref{sec:hdp} describes the $\hDP$ tetrahedron.
Section~\ref{sec:buffer} proves that baryons require this skeleton
while mesons do not.  Section~\ref{sec:occupation} analyzes the
3-of-4 occupation rule and the open fourth vertex.
Section~\ref{sec:color} establishes color as vertex identity and
derives the color selection rules.  Section~\ref{sec:su3} identifies
SU(3) as emergent bookkeeping.  Section~\ref{sec:gluons} states the
gluon count correspondence and defers its derivation to
companion~16~\cite{c16}.

%======================================================================
\section{The hDP Tetrahedral Skeleton}
\label{sec:hdp}
%======================================================================

\subsection{Hybrid Dipole Particles}

A hybrid Dipole Particle ($\hDP$) is a DP that combines both
electromagnetic ($\eCP$) and strong-force ($\qCP$) character at its
two poles.  Unlike the purely electromagnetic $\eDP$ (which carries
charge $\pm e$) or the purely strong $\qDP$ (which carries charge
$\pm e_q$), the $\hDP$ carries a mixed charge at each pole --- one
pole dominated by $\eCP$ character, the other by $\qCP$ character.

This mixed character makes the $\hDP$ a mediating entity: it can
couple to both the electromagnetic and strong SSV fields, allowing it
to buffer interactions between quarks of like charge.

\subsection{Spontaneous Nucleation of hDP Tetrahedra}

The Dipole Sea contains $\eDP$, $\qDP$, and $\hDP$ populations in
thermal equilibrium.  The $\hDP$ subpopulation spontaneously nucleates
tetrahedral structures for the following reasons:

\begin{enumerate}
\item \textbf{Alternating charge geometry.}  A tetrahedron with
  alternating $+\hDP$ and $-\hDP$ at its four vertices (two positive,
  two negative) minimizes electrostatic energy relative to any other
  four-vertex configuration.  The alternating geometry produces an
  attractive potential well that stabilizes the tetrahedron against
  thermal disruption.

\item \textbf{ZBW resonance locking.}  Each $\hDP$ oscillates at
  the Zitterbewegung frequency (companion~4~\cite{c4}).  When four
  $\hDP$s are within one lattice spacing, their ZBW oscillations
  can phase-lock into a collective resonance mode.  This resonance
  contributes additional binding energy, setting the equilibrium
  size of the tetrahedron.

\item \textbf{600-cell GP site availability.}  The four vertices of
  the $\hDP$ tetrahedron must correspond to Grid Points (GPs) on
  the 600-cell lattice (Version~16~\cite{v16}).  The 600-cell's
  tetrahedral subgroup ($T_d$, a subgroup of $H_4$) provides exactly
  the GP configurations required.  The lattice constrains which
  tetrahedra can form and enforces the discrete geometry.
\end{enumerate}

\begin{remark}[Stability regime]
The $\hDP$ tetrahedron is more stable than a pure $\eDP$ tetrahedron
(which lacks strong-force binding) and less stable than a pure $\qDP$
tetrahedron (which lacks electromagnetic mediation).  This intermediate
stability is precisely what makes the $\hDP$ tetrahedron a suitable
quark bonding scaffold: it is stable enough to persist on the timescale
of hadron dynamics, but flexible enough to accommodate quark bonding
transitions.
\end{remark}

\begin{remark}[Cosmological context]
The $\hDP$ tetrahedron population may be in dynamic thermal
equilibrium at the current temperature of the universe, or may have
frozen out during cosmic cooling as a preferred configuration.  Both
framings are consistent with CPP; the cosmological history is an open
problem (§\ref{sec:open}).
\end{remark}

\subsection{Geometry of the hDP Tetrahedron}

A regular tetrahedron has four vertices, six edges, and four faces.
The $T_d$ point symmetry group of the regular tetrahedron has order
24, containing:
\begin{itemize}[noitemsep]
\item 1 identity element
\item 8 proper rotations ($C_3$ and $C_3^2$ about each of 4 axes)
\item 3 $C_2$ rotations
\item 6 $S_4$ improper rotations
\item 6 $\sigma_d$ reflection planes
\end{itemize}

The $C_3$ rotational axis through any vertex and the centroid of the
opposite face permutes the three remaining vertices cyclically.  This
3-fold symmetry is the geometric origin of the color triplet, as
established in §\ref{sec:color}.

%======================================================================
\section{Why Baryons Require the hDP Skeleton}
\label{sec:buffer}
%======================================================================

\subsection{The Charge Buffer Theorem}

\begin{theorem}[Baryons require the hDP tetrahedral skeleton]
\label{thm:buffer}
Three like-sign quarks cannot form a stable bound state without
the $\hDP$ tetrahedral skeleton as a charge buffer.
\end{theorem}

\begin{proof}
Consider three quarks of the same sign charge (e.g., three $+\qCP$
quarks, as in the two up quarks and one down quark of a proton where
both up quarks carry charge $+2e/3$).  Without the $\hDP$ tetrahedron:

The SSV field of each $+\qCP$ quark exerts a repulsive force on every
other $+\qCP$ quark at distance $r$:
\begin{equation}
F_{\rm rep}(r) = \frac{\as\,\hbar c}{r^2} > 0
\quad \text{(repulsive, same-sign quarks)}.
\label{eq:repulsion}
\end{equation}

For the three quarks to remain in proximity, a binding energy
exceeding $F_{\rm rep} \cdot \Delta r$ must exist.  The only
available binding mechanism in CPP at sub-fm scales is:
(a) direct qDP chain attraction between opposite-sign quarks, or
(b) mediated attraction via an intervening opposite-sign CP.

Two up quarks ($+2e/3$ each) have no direct qDP chain attraction
between them.  Mechanism (b) requires an opposite-sign mediator
at each vertex.  The $\hDP$ tetrahedron provides exactly this:
each vertex presents an opposite-sign $\hDP$ pole to the incoming
quark, creating a local attraction that overcomes the inter-quark
repulsion.

Without the $\hDP$ tetrahedron, the repulsive force
Eq.~\eqref{eq:repulsion} dominates at all scales below $r_0$, and
the three-quark configuration dissociates on a timescale
$\tau \sim \hbar / F_{\rm rep} \cdot l_P \ll t_P$.  The bound
state cannot exist.

With the $\hDP$ tetrahedron, each quark bonds to an opposite-sign
$\hDP$ vertex.  The quark-vertex attraction exceeds the inter-quark
repulsion, stabilizing the three-quark configuration.
\end{proof}

\subsection{Why Mesons Do Not Require the Skeleton}

\begin{proposition}[Mesons bind without the hDP skeleton]
\label{prop:meson_no_skeleton}
A quark-antiquark ($q\bar{q}$) pair does not require the $\hDP$
tetrahedral skeleton for binding.
\end{proposition}

\begin{proof}
A quark ($+\qCP$) and antiquark ($-\qCP$) carry opposite sign charges.
They attract directly via the qDP chain mechanism of companion~14~\cite{c14}:
\begin{equation}
V_{\rm meson}(r) = -\frac{\as\,\hbar c}{r} + \sigma\,r.
\label{eq:meson_potential}
\end{equation}
No charge buffering is required.  The qDP chain provides sufficient
binding at all distances from $r = 0$ to string-breaking.
\end{proof}

\begin{remark}[Why a meson does not acquire a tetrahedral skeleton]
If a tetrahedral $\hDP$ skeleton were drawn into the meson's qDP chain
region, its open vertices would immediately attract additional quarks
from the Dipole Sea.  This would convert the meson into a baryon-like
configuration with three quark bonding sites.  The meson's two-body
symmetry would be broken.  The resulting configuration is either a
higher-energy state (which relaxes back to the plain meson) or decays
into separate hadrons.  The meson therefore actively resists tetrahedral
skeleton formation as a matter of energy minimization.
\end{remark}

%======================================================================
\section{The 3-of-4 Occupation Rule and Net Baryon Charge}
\label{sec:occupation}
%======================================================================

\subsection{Why Three Vertices Are Occupied}

The $\hDP$ tetrahedron has four vertices.  Stable baryons occupy
exactly three.  This is the \emph{3-of-4 occupation rule}.

\begin{proposition}[3-of-4 occupation rule]
\label{prop:three_of_four}
In a stable baryon, exactly three of the four $\hDP$ tetrahedral
vertices are occupied by quarks.  Full four-vertex occupation is
unstable.
\end{proposition}

\begin{proof}[Mechanistic argument]
Consider the energy balance at the fourth vertex when three vertices
are already occupied.

\textbf{Symmetry equalization:} When all four vertices are occupied,
the $T_d$ symmetry of the tetrahedron is fully realized.  This
equalization distributes binding energy equally across all four
bonds, reducing the energy advantage of each individual bond.  Each
bond is weaker than the bonds in the three-vertex case, where the
asymmetry of the occupied-vs-open configuration concentrates binding
energy into the three occupied vertices.

\textbf{SSV invasion:} The two like-sign quarks at vertices $V_1$
and $V_2$ (e.g., two $+2e/3$ up quarks) each emit repulsive SSV
fields.  At the fourth vertex $V_4$, these repulsive SSV fields
from $V_1$ and $V_2$ partially cancel the attractive SSV of the
$\hDP$ vertex charge.  The effective binding energy at $V_4$ is:
\begin{equation}
E_{\rm bind}(V_4) = E_{\rm vertex}^{\rm attract}
- \sum_{i \in \{V_1,\, V_2\}} \frac{\as\,\hbar c}{|V_i - V_4|}
< E_{\rm bind}(V_{1,2,3}).
\label{eq:V4_reduced}
\end{equation}
When $E_{\rm bind}(V_4)$ falls below the thermal disruption threshold
$k_B T_{\rm QCD} \approx \Lambda_{\rm QCD} \approx 213$~MeV, the
fourth quark is thermally ejected.

The combined effect of symmetry equalization and SSV invasion reduces
the fourth vertex binding energy below stability, leaving it
perpetually open in the ground state.
\end{proof}

\subsection{Net Baryon Charge from the Open Vertex}

The open fourth vertex is not electroneutral.  It exposes a net
charge from the unshielded $\hDP$ vertex pole.

\begin{proposition}[Net baryon charge]
\label{prop:baryon_charge}
The net charge of a baryon equals the charge of the unoccupied
fourth vertex of the $\hDP$ tetrahedral skeleton.
\end{proposition}

\begin{proof}
Each occupied vertex is charge-shielded: the quark bonded to it
provides an opposite-sign CP that neutralizes the $\hDP$ vertex
charge locally.  The three occupied vertices contribute zero net
charge to the baryon's long-range field.

The unoccupied fourth vertex exposes its $\hDP$ pole charge without
shielding.  For the proton:
\begin{itemize}[noitemsep]
\item Two up quarks ($+2e/3$ each) bond to two $-\hDP$ vertices:
  net charge at each occupied vertex $= +2e/3 - 2e/3 = 0$.
\item One down quark ($-e/3$) bonds to one $+\hDP$ vertex:
  net charge at that vertex $= -e/3 + e/3 = 0$.
\item The unoccupied fourth vertex exposes charge $+e/3$.
\end{itemize}
Summing: $0 + 0 + 0 + e/3 = +e/3$ from the $\hDP$ vertex, plus
the quark charges $+2e/3 + 2e/3 - e/3 = +e$.  The total is:
\begin{equation}
Q_{\rm proton} = (+2e/3) + (+2e/3) + (-e/3) = +e.
\label{eq:proton_charge}
\end{equation}
This matches the observed proton charge exactly~\cite{pdg2022}.
\end{proof}

\begin{remark}[Neutron charge]
For the neutron ($udd$): two down quarks ($-e/3$ each) bond to two
$+\hDP$ vertices; one up quark ($+2e/3$) bonds to one $-\hDP$
vertex.  The fourth vertex exposes $-e/3$ from the unoccupied
$\hDP$ pole.  Net charge:
\begin{equation}
Q_{\rm neutron} = (+2e/3) + (-e/3) + (-e/3) = 0.
\label{eq:neutron_charge}
\end{equation}
The unoccupied $-e/3$ vertex is exactly cancelled by the asymmetry
in the quark configuration, giving the neutron its observed zero
charge~\cite{pdg2022}.
\end{remark}

%======================================================================
\section{Color Charge as Vertex Identity}
\label{sec:color}
%======================================================================

\subsection{Definition}

\begin{definition}[Color charge as vertex label]
\label{def:color}
The color charge of a quark in a baryon is its vertex identity:
the label (red, green, or blue) assigned to the specific $\hDP$
tetrahedral vertex to which it is bonded.  Color is not an intrinsic
quantum number carried by the quark; it is a relational property
describing the quark's position in the tetrahedral bonding structure.
\end{definition}

The three color labels correspond to the three $C_3$-related base
vertices of the tetrahedron.  Under the $C_3$ rotation axis through
the apex vertex and the centroid of the opposite face, the three
base vertices permute cyclically:
\begin{equation}
C_3: \quad \text{red} \to \text{green} \to \text{blue} \to \text{red}.
\label{eq:C3_permutation}
\end{equation}
This cyclic permutation is the color rotation symmetry.

\subsection{Color Selection Rules Derived from Vertex Geometry}

\begin{theorem}[Color neutrality]
\label{thm:color_neutral}
A color-neutral hadron is one in which all three quark-bonding
vertices of the $\hDP$ tetrahedron are occupied, with no net
unshielded vertex charge.
\end{theorem}

\begin{proof}
By Definition~\ref{def:color}, each color label corresponds to one
tetrahedral vertex.  ``Color neutral'' in QCD means the sum of color
charges is zero.  In CPP, this corresponds to complete occupation of
all three quark-bonding vertices: red $+$ green $+$ blue $= $ all
three vertices occupied $=$ no net unshielded $\hDP$ vertex SSV
field extending beyond the hadron radius.  The three quarks'
opposite-sign bonds to the three $\hDP$ vertices neutralize all
three vertex charges, leaving only the open fourth vertex exposed
as the net baryon charge (§\ref{sec:occupation}).
\end{proof}

\begin{theorem}[Prohibition on free quarks]
\label{thm:no_free_quark}
A single quark cannot exist as a free particle.
\end{theorem}

\begin{proof}
By the confinement theorem of companion~14 (Theorem~1 in~\cite{c14}), any
isolated quark requires a qDP chain extending to infinity, which
costs infinite energy.  In the tetrahedral skeleton picture, a
single quark bonded to one vertex would expose three unshielded
$\hDP$ vertex charges (the three unoccupied vertices), producing
a strong long-range SSV field that attracts additional quarks
from the Dipole Sea until the three-vertex occupation is achieved.
Both the chain topology (companion~14) and the vertex geometry
(this paper) independently prohibit free quarks.
\end{proof}

\begin{proposition}[Allowed color configurations]
\label{prop:allowed_configs}
The only stable color configurations are:
\begin{enumerate}[noitemsep]
\item Three quarks of distinct vertex labels ($qqq$): baryon.
\item Three antiquarks of distinct vertex labels ($\bar{q}\bar{q}\bar{q}$):
  antibaryon.
\item One quark and one antiquark via direct qDP chain ($q\bar{q}$):
  meson (no tetrahedral skeleton required).
\end{enumerate}
\end{proposition}

\begin{proof}[Proof sketch]
Configuration 1 follows from Theorem~\ref{thm:color_neutral}: the
only way to achieve color neutrality on the $\hDP$ tetrahedron is to
occupy all three quark-bonding vertices with quarks of distinct labels.
Configuration 2 is the antibaryon analog by charge conjugation.
Configuration 3 follows from Proposition~\ref{prop:meson_no_skeleton}:
the meson binds without the tetrahedral skeleton, so no vertex-label
constraint applies; color neutrality is achieved by quark-antiquark
color-anticolor pairing via the qDP chain.

No other stable configuration exists because:
\begin{itemize}[noitemsep]
\item Two quarks on the tetrahedron leave one vertex open with
  unshielded charge, attracting a third quark.
\item Four quarks are unstable by Proposition~\ref{prop:three_of_four}.
\item Exotic configurations ($qqqqq$, tetraquarks) require double-skeleton
  structures — possible as transient resonances, deferred to
  Open Problem~\ref{op:exotic}.
\end{itemize}
\end{proof}

\subsection{Color Change as Vertex Transition}

A color change --- a quark transitioning from one color state to
another --- corresponds physically to a quark changing its bonding
vertex on the $\hDP$ tetrahedron.  This is a real mechanical
transition mediated by a $\qDP$ chain oscillation (a gluon in the
QCD language):

\begin{equation}
q_{\rm red} + g_{r\bar{g}} \to q_{\rm green}
\quad \Longleftrightarrow \quad
q_{V_1} + \text{(qDP transition)} \to q_{V_2}.
\label{eq:color_change}
\end{equation}

The gluon is not a particle carrying abstract color-anticolor quantum
numbers; it is a $\qDP$ chain excitation whose mode structure enables
the quark to detach from vertex $V_1$ and attach to vertex $V_2$.
The full mode-to-generator correspondence (which vertex transitions
correspond to which Gell-Mann generators) is established in
companion~16~\cite{c16}.

%======================================================================
\section{SU(3) as Emergent Bookkeeping}
\label{sec:su3}
%======================================================================

\begin{proposition}[SU(3) as permutation algebra]
\label{prop:su3_bookkeeping}
The SU(3) color algebra is the mathematical description of the
permutation symmetry of 3-of-4 vertex occupation on the $\hDP$
tetrahedron.  It is not a fundamental gauge symmetry; it is the
correct algebra for bookkeeping the observed color selection rules.
\end{proposition}

The argument proceeds in three steps.

\textbf{Step 1 — Z$_3$ center.}
The $C_3$ rotation of Eq.~\eqref{eq:C3_permutation} generates the
cyclic group $\mathbb{Z}_3$.  In the representation theory of SU(3),
the center of the group is exactly $\mathbb{Z}_3$: color charge is a
$\mathbb{Z}_3$ quantum number.  CPP derives this as the discrete
rotational degeneracy of the tetrahedral $C_3$ axis, not as a
postulate.

\textbf{Step 2 — Triplet and antitriplet representations.}
The three occupied vertices of the tetrahedron transform as the
fundamental \textbf{3} representation of SU(3) under color rotation
($C_3$ permutation).  The three antiquark vertices transform as the
conjugate $\bar{\textbf{3}}$.  Color neutrality requires either
$\mathbf{3} \otimes \bar{\mathbf{3}}$ (meson) or
$\mathbf{3} \otimes \mathbf{3} \otimes \mathbf{3}$ (baryon) to
contain the singlet \textbf{1}.  Both do:
\begin{align}
\mathbf{3} \otimes \bar{\mathbf{3}} &= \mathbf{1} \oplus \mathbf{8},
\label{eq:meson_decomp}\\
\mathbf{3} \otimes \mathbf{3} \otimes \mathbf{3} &=
\mathbf{1} \oplus \mathbf{8} \oplus \mathbf{8} \oplus \mathbf{10}.
\label{eq:baryon_decomp}
\end{align}
These are standard SU(3) Clebsch-Gordan decompositions.  In CPP, they
express how vertex-occupation configurations combine to give
color-neutral states.

\textbf{Step 3 — SU(3) as the minimal algebra fitting the geometry.}
Any algebra that correctly describes permutations of three
distinguishable objects subject to the $\mathbb{Z}_3$ cyclic
constraint and the constraint that only neutral combinations are
stable must contain SU(3) as a substructure.  SU(3) is the
\emph{simplest} Lie algebra with center $\mathbb{Z}_3$ and a faithful
3-dimensional representation.  The 600-cell tetrahedral geometry
selects this algebra automatically, without requiring it to be
postulated.

\begin{remark}[What SU(3) does and does not explain]
SU(3) correctly predicts all color selection rules, the hadron
multiplet structure, and the number of gluons.  What it does not
explain is \emph{why} color charge exists, \emph{why} only three
colors, and \emph{why} baryons are color-neutral.  CPP answers these
questions mechanically: color exists because quarks bond to $\hDP$
tetrahedral vertices; there are three colors because the tetrahedron
has three quark-bonding vertices; baryons are color-neutral because
three-vertex occupation shields all vertex charges.  SU(3) is the
correct algebra for this geometry, but the geometry is more
fundamental than the algebra.
\end{remark}

%======================================================================
\section{The Gluon Octet: Count and Deferral}
\label{sec:gluons}
%======================================================================

The 8-dimensional adjoint representation of SU(3) --- the gluon octet
--- corresponds in CPP to the 8 distinct $\qDP$ chain modes that mediate
vertex transitions on the $\hDP$ tetrahedron.

The count $8 = 1 + 3 + 4$ from companion~14~\cite{c14} (one central
chain, three middle-layer chains at 120°, four outer-layer chains at 60°)
matches the 8 Gell-Mann generators exactly.  In the SU(3) algebra,
these 8 generators decompose as:
\begin{itemize}[noitemsep]
\item 6 off-diagonal generators: color-changing transitions
  (quark moves from one vertex to another).
\item 2 diagonal generators (the Cartan subalgebra $T_3$, $Y$):
  color-preserving transitions that change energy state without
  changing vertex.
\end{itemize}

The explicit correspondence between the 1+3+4 chain layer modes and
the 8 Gell-Mann matrices --- identifying which chain mode produces
which vertex transition and mapping it to the appropriate generator
--- requires working out the ZBW oscillation modes of the $\hDP$
tetrahedral skeleton in detail.  This is a non-trivial calculation
deferred to companion~16~\cite{c16}.

\begin{remark}[Why 8 and not 9]
SU(3) has 8 generators rather than 9 because one linear combination
of the diagonal generators corresponds to the identity transformation
(the U(1) factor that is gauged away in QCD).  In CPP, this
corresponds to the one chain mode that produces no net vertex
transition and no net color change --- a ``breathing mode'' of the
entire $\hDP$ tetrahedron.  This mode is suppressed because the
$\hDP$ tetrahedron's alternating charge geometry forbids uniform
radial breathing (it would require all four vertices to move
simultaneously, violating the 3-of-4 occupation constraint).
\end{remark}

%======================================================================
\section{Particle Structure and the Cage Hierarchy}
\label{sec:cages}
%======================================================================

The $\hDP$ tetrahedron of the baryon is the simplest instance of a
broader pattern in CPP: all massive particles nucleate around
successive geometric cages of DPs on 600-cell shell structures.
The quark and lepton mass hierarchy reflects the cage hierarchy:

\begin{table}[h]
\centering
\caption{Quark and lepton cage structure in CPP.  Mass increases
with cage complexity.  The $\hDP$ tetrahedron relevant to this
paper is the baryon bonding scaffold, distinct from (but related
to) the individual quark cage structures.}
\label{tab:cages}
\renewcommand{\arraystretch}{1.3}
\begin{tabular}{@{}llll@{}}
\toprule
Particle & Central CP & Outer cage & Approx.\ mass \\
\midrule
Electron & $\eCP$ & None (bare) & 0.511 MeV \\
Up/Down quark & $\qCP$ & None / single $\hDP$ & $\sim$2--5 MeV \\
Muon & $\eCP$ & Tetrahedral ($\hDP$, 4 vertices) & 105.7 MeV \\
Strange quark & $\qCP$ & Tetrahedral ($\hDP$, 4 vertices) & $\sim$95 MeV \\
Tau & $\eCP$ & Icosahedral (12 vertices) & 1777 MeV \\
Charm quark & $\qCP$ & Icosahedral (12 vertices) & $\sim$1275 MeV \\
Bottom quark & $\qCP$ & Dodecahedral (20 vertices) & $\sim$4180 MeV \\
Top quark & $\qCP$ & Fullerene-like ($\sim$60 vertices) & $\sim$173 GeV \\
\midrule
W boson & None & 6-$\hDP$ loop (ribbon) & 80.4 GeV \\
Z boson & None & Icosahedral ($\hDP$, 12 vertices) & 91.2 GeV \\
Higgs & None & Dodecahedral ($\hDP$, 20 vertices) & 125.1 GeV \\
\bottomrule
\end{tabular}
\end{table}

The generational hierarchy of quarks and leptons emerges from this
cage progression: first generation (up/down, electron) has minimal
or no cage; second generation (strange/charm, muon) has a
tetrahedral cage; third generation (bottom/top, tau) has icosahedral,
dodecahedral, or fullerene cages.  Mass is cage energy --- the ZBW
oscillation energy of the DP cage plus the central CP's Compton
frequency contribution (companion~4~\cite{c4}).

The bosons (W, Z, Higgs) have no central unpaired CP and therefore
no intrinsic charge.  They are pure $\hDP$ cage structures whose
geometry determines their mass and interaction cross-section.

\begin{remark}[Scope of this section]
Table~\ref{tab:cages} is a structural overview establishing the
context of the $\hDP$ tetrahedron within the broader CPP particle
structure.  The derivation of individual particle masses from cage
geometry is the subject of a future companion in the SM series.
\end{remark}

%======================================================================
\section{Consistency with Previous Companions}
\label{sec:consistency}
%======================================================================

\begin{itemize}[noitemsep]
\item The \emph{qDP chaining mechanism} and \emph{Cornell potential}
  of companion~14~\cite{c14} are unchanged.  Color is a layer above
  the chain mechanism, not in conflict with it.  The confinement
  theorem (infinite energy for free quarks) is complemented by the
  vertex-geometry prohibition of §\ref{sec:color}.
\item The \emph{ZBW mass mechanism} of companion~4~\cite{c4} applies
  to quarks: each quark's mass arises from its Compton-scale ZBW
  oscillation, independent of its color/vertex identity.
\item The \emph{fractional charge} of quarks ($e/3$, $2e/3$) derived
  in companion~2~\cite{c2} from cage projection is independent of
  the color/vertex assignment.  Charge fraction is set by the
  individual quark cage geometry; color is set by which
  $\hDP$ tetrahedral vertex the quark occupies.
\item The \emph{CP Exclusion Rule} of companion~1~\cite{c1} governs
  the minimum PSR of each quark CP, setting the scale at which
  quark-vertex bonds can form and break.
\item The \emph{600-cell GP structure} of Version~16~\cite{v16}
  provides the lattice sites for both the individual quark cages
  and the $\hDP$ tetrahedral skeleton.  The $T_d$ tetrahedral
  subgroup of $H_4$ is the geometric foundation of both.
\end{itemize}

No new postulates beyond those of Version~16 are required.

%======================================================================
\section{Open Problems}
\label{sec:open}
%======================================================================

\begin{enumerate}
\item \label{op:hdp_size}
  \textbf{First-principles derivation of the hDP tetrahedron size.}
  The equilibrium size of the $\hDP$ tetrahedron is governed by the
  balance between electrostatic attraction of the alternating charge
  geometry and ZBW resonance locking.  A quantitative calculation
  requires the $\hDP$ ZBW frequency and the 600-cell GP spacing at
  the relevant shell --- an extension of the ZBW mass calculation
  of companion~4~\cite{c4} to the hybrid DP sector.

\item \textbf{Cosmological freezing of the hDP tetrahedron population.}
  Whether the $\hDP$ tetrahedron population is in current thermal
  equilibrium or froze out during cosmic cooling determines the
  baryon density of the universe.  A CPP treatment of this
  question requires the full thermodynamics of the DP Sea at
  temperatures $T \sim \Lambda_{\rm QCD}$, connecting to the
  Big Bang and baryon asymmetry.

\item \label{op:fourth_vertex}
  \textbf{Rigorous proof of the 3-of-4 occupation rule.}
  Proposition~\ref{prop:three_of_four} gives a mechanistic argument.
  A rigorous proof requires computing $E_{\rm bind}(V_4)$ from
  the $\hDP$ vertex SSV field and the like-charge quark SSV invasion
  explicitly, showing the result is below the QCD thermal threshold
  $k_B T_{\rm QCD}$ for all stable baryon configurations.

\item \label{op:exotic}
  \textbf{Exotic hadrons from double-skeleton structures.}
  Tetraquarks ($qq\bar{q}\bar{q}$) and pentaquarks ($qqqq\bar{q}$)
  have been observed at the LHC~\cite{pdg2022}.  In CPP these
  may correspond to double-$\hDP$ skeleton configurations.
  The stability conditions and mass spectrum of such structures
  are an open problem.

\item \textbf{Gluon mode-to-generator correspondence.}
  The explicit mapping between the $1 + 3 + 4 = 8$ qDP chain
  layer modes and the 8 Gell-Mann generators, including the
  identification of the 6 color-changing and 2 diagonal modes,
  is deferred to companion~16~\cite{c16}.
\end{enumerate}

%======================================================================
\section{Conclusion}
\label{sec:conclusion}
%======================================================================

\begin{enumerate}
\item \textbf{hDP tetrahedron} (§\ref{sec:hdp}): The Dipole Sea
  spontaneously nucleates $\hDP$ tetrahedral skeletons stabilized
  by alternating charge geometry, ZBW resonance locking, and
  600-cell GP site availability.

\item \textbf{Charge buffer theorem} (Theorem~\ref{thm:buffer}):
  Three like-charge quarks cannot form a stable bound state without
  the $\hDP$ tetrahedron as a charge buffer.  Mesons do not require
  the skeleton because quark-antiquark pairs attract directly via
  the qDP chain.

\item \textbf{3-of-4 occupation} (Proposition~\ref{prop:three_of_four}):
  Exactly three tetrahedral vertices are occupied in a stable baryon.
  The fourth vertex is destabilized by symmetry equalization and
  SSV invasion from like-charge quarks.

\item \textbf{Net baryon charge} (Proposition~\ref{prop:baryon_charge}):
  The net charge of a baryon is the unshielded charge of the open
  fourth vertex.  Proton charge $+e$ and neutron charge $0$ are
  derived without additional assumptions.

\item \textbf{Color as vertex identity} (Definition~\ref{def:color}):
  Color charge is not an intrinsic quantum number; it is a quark's
  vertex label on the $\hDP$ tetrahedral skeleton.  Red, green, and
  blue are names for three $C_3$-related tetrahedral vertices.

\item \textbf{Color selection rules} (Theorems~\ref{thm:color_neutral},
  \ref{thm:no_free_quark}, Proposition~\ref{prop:allowed_configs}):
  Color neutrality, the prohibition on free quarks, and the allowed
  hadron configurations ($q\bar{q}$, $qqq$, $\bar{q}\bar{q}\bar{q}$)
  all follow from charge geometry without gauge postulates.

\item \textbf{SU(3) as bookkeeping} (Proposition~\ref{prop:su3_bookkeeping}):
  SU(3) is the permutation algebra of 3-of-4 vertex occupation.
  It is the correct and useful algebra for the geometry, but it is
  derived from the geometry --- not the other way around.

\item \textbf{Gluon count} (§\ref{sec:gluons}): The 8 gluon modes
  correspond to the $1 + 3 + 4 = 8$ qDP chain layer modes of
  companion~14.  The explicit mode-to-generator mapping is deferred
  to companion~16.
\end{enumerate}

Color charge, in CPP, is not a mystery requiring a new gauge
symmetry.  It is the mechanical consequence of quarks bonding to
the vertices of a geometric skeleton that the Dipole Sea
spontaneously provides.  The skeleton is stable, its vertices are
distinguishable, and the rules of color follow from the rules of
which vertices can be occupied.

%======================================================================
\section*{Acknowledgments}
%======================================================================

The conception of the $\hDP$ tetrahedral skeleton as the physical
substrate of baryon structure, and the identification of color charge
as an artifact of the resonant bonding modes of quarks on the
tetrahedral skeleton, are the intellectual contribution of Thomas Lee
Abshier, ND, developed over multiple years of theoretical development.
The formalization of color as vertex identity, the charge buffer
theorem, and the 3-of-4 occupation argument were developed in
collaboration with Claude (Anthropic) on 19~March 2026.

%======================================================================
\begin{thebibliography}{99}

\bibitem{v16}
T.~L. Abshier,
``Mechanistic Derivation of Relativistic Effects via Space Stress
Vector (SSV) in the Dipole Sea,''
Version~16, 2026 (main paper).

\bibitem{c1}
T.~L. Abshier,
``The Absolute Moment Postulate,''
Version~3, 2026 (companion~1).

\bibitem{c2}
T.~L. Abshier,
``Microscopic Origin of the Dipole Sea Stiffness~$C$,''
Version~2, 2026 (companion~2).

\bibitem{c4}
T.~L. Abshier,
``Inertial Mass from Zitterbewegung,''
Version~2, 2026 (companion~4).

\bibitem{c14}
T.~L. Abshier,
``Quark Confinement and the Cornell Potential from qDP Chaining
in Conscious Point Physics,''
Version~1.1, 2026 (companion~14).

\bibitem{c16}
T.~L. Abshier,
``Gluon Modes and the SU(3) Generator Correspondence in CPP,''
2026 (companion~16, in preparation).

\bibitem{gross1973}
D.~J. Gross and F.~Wilczek,
``Ultraviolet Behavior of Non-Abelian Gauge Theories,''
\textit{Phys.\ Rev.\ Lett.}, \textbf{30}, 1343--1346, 1973.

\bibitem{politzer1973}
H.~D. Politzer,
``Reliable Perturbative Results for Strong Interactions?''
\textit{Phys.\ Rev.\ Lett.}, \textbf{30}, 1346--1349, 1973.

\bibitem{pdg2022}
R.~L. Workman \textit{et al.}\ (Particle Data Group),
``Review of Particle Physics,''
\textit{Prog.\ Theor.\ Exp.\ Phys.}, \textbf{2022}, 083C01, 2022.

\end{thebibliography}

\end{document}
